% Generated draft: requires mathematical and source review.
\documentclass[11pt]{article}
\usepackage[margin=25mm]{geometry}
\usepackage{fontspec}
\setmainfont{Latin Modern Roman}
\usepackage{amsmath,amssymb,mathtools,unicode-math}
\setmathfont{Latin Modern Math}
\usepackage{enumitem,needspace,xurl,hyperref}
\setlength{\parindent}{0pt}
\setlength{\parskip}{5pt}
\setlength{\emergencystretch}{3em}
\newcommand{\N}{\mathbb{N}}
\newcommand{\Z}{\mathbb{Z}}
\newcommand{\Q}{\mathbb{Q}}
\newcommand{\R}{\mathbb{R}}
\newcommand{\C}{\mathbb{C}}
\newcommand{\F}{\mathbb{F}}
\newcommand{\A}{\mathbb{A}}
\newcommand{\PP}{\mathbb{P}}
\newcommand{\E}{\mathbb{E}}
\newcommand{\eps}{\varepsilon}
\newcommand{\dd}{\,\mathrm{d}}
\begin{document}
\section*{UnsolvedMath ID: 30004461}
\textbf{Source problem code:} catalogue entry
\subsection{Definitions and mathematical statement}
Let \(\mathcal C\) denote a category of bornological algebras on which the intended analytic and periodic cyclic homology theories are both defined. A bornology specifies the bounded subsets of an algebra; its compatibility with the algebra operations is understood in the chosen bornological category. The supplied record does not specify the ground field, completeness assumptions, or precise bornological category, so these remain parameters of the formulation.

For \(A\in\mathcal C\), write \(\operatorname{HA}_*(A)\) and \(\operatorname{HP}_*(A)\) for the intended analytic and periodic cyclic homology, respectively, with compatible coefficient and grading conventions. Here these symbols name the theories referred to in the source; their precise constructions are not supplied.

\textbf{Comparison problem.} Determine conditions on \(A\in\mathcal C\) under which there is an isomorphism of graded homology objects
\[
\operatorname{HA}_*(A)\cong\operatorname{HP}_*(A).
\]
Equivalently, after fixing the preceding conventions, determine the class
\[
\mathcal C_{\mathrm{eq}}
=
\bigl\{A\in\mathcal C:\operatorname{HA}_*(A)\cong\operatorname{HP}_*(A)\bigr\}.
\]
In particular, can the comparison results that the supplied record attributes to the analytically quasi-free setting be extended to a broader class of bornological algebras?

\textbf{Unresolved qualifications.} The record does not define analytically quasi-free algebras or state the hypotheses and conclusions of the comparison results it invokes. It also does not specify a comparison morphism, the direction of such a morphism, or whether the requested isomorphism must be natural. Consequently, the displayed condition records the comparison target without asserting a canonical comparison map or an isomorphism for every bornological algebra. A fully specified theorem statement requires these missing conventions and hypotheses.
\subsection{Short English statement}
Determine when analytic cyclic homology and periodic cyclic homology agree for a bornological algebra, and whether the comparison regime attributed by the source to analytically quasi-free algebras extends further. The category, theory conventions, and precise existing comparison results remain unspecified. The supplied open status is a dated curation claim.
\subsection{Sources}
\begin{itemize}
\item Local UnsolvedMath registry: supplied source record for permanent ID \(30004461\), titled \emph{Comparison of Analytic and Periodic Cyclic Homology}, in the Oberwolfach Reports Open Problems collection. Exact source problem code: \texttt{catalogue entry}. No individual proposer is identified; the supplied proposed year is \(2020\).
\item Supplied report citation: \emph{Non-Commutative Geometry and Cyclic Homology} (2021), Oberwolfach Reports, pp.~919--960, DOI \texttt{10.4171/owr/2020/17}. Original OWR report: \url{https://doi.org/10.4171/owr/2020/17}.
\item Supplied literature reference labeled Meyer--Mukherjee: \url{https://arxiv.org/abs/2109.01470}. The supplied record gives no title or specific theorem reference.
\item Literature assessment checked at \texttt{2026-08-22}: the supplied curation labels the problem open and reports that analytic cyclic homology has strong formal properties and agrees with periodic cyclic homology in special settings, while no general comparison theorem settles the extracted question. This is supplied dated curation, not an independently verified assessment.
\item The supplied record labels its statement \texttt{corrected\_verified} and says that the official report was checked and the analytically quasi-free regime included. These are source-verification claims of the supplied record; no independent source visit or verification is asserted here.
\end{itemize}
\end{document}
